%%%%%%%%%%%%%%%%%%%%%%%%%%%%%%%%%%%%%%%%%%%%%%%%%%%%%%%%%%%%%%%%%%%%%%
% Lớp 12 | Ứng dụng đạo hàm để khảo sát hàm số | Cực trị của hàm số
% Dạng 2D1?2-5: Tìm 𝑚 để hàm số, đồ thị hàm số trùng phương có cực trị thỏa mãn điều
% Số câu: 2  (sinh tự động bởi tach_cau.py; từ nay thầy sửa trực tiếp file này)
%%%%%%%%%%%%%%%%%%%%%%%%%%%%%%%%%%%%%%%%%%%%%%%%%%%%%%%%%%%%%%%%%%%%%%

%%%%%%%%%%%%%%%% TRẮC NGHIỆM 4 PHƯƠNG ÁN %%%%%%%%%%%%%%%%

% nguồn: Phan1_Lop_12_Dai_so_Chuong_1
\begin{ex}%[2D1V2-5]%[ma:93a6e0c4]
	Cho hàm số $y=(m+1)x^4-mx^2+3$. Tìm tất cả các giá trị thực của tham  số $m$ để hàm số có ba điểm cực trị.
	\choice 
		{$m\in (-\infty;-1)\cup [0;+\infty)$} 
		 {$m\in (-1;0)$}
		{$m\in (-\infty;-1]\cup  [0;+\infty)$}
		{\True $m\in (-\infty;-1)\cup  (0;+\infty)$}
	\loigiai{ 
		 $y^\prime= 4(m+1)x^3-2mx=2x[2(m+1)x^2-m]$\\
		$y^\prime =0\Leftrightarrow \hoac{&x=0\\&2(m+1)x^2-m=0\,\,(*)}  $\\
		Hàm số có ba điểm cực trị khi phương trình $(*)$ có hai nghiệm phân biệt khác  $0$
		 $$\Leftrightarrow \dfrac{m}{m+1}>0\Leftrightarrow \hoac{&m<-1\\&m>0.} $$
	 }
\end{ex}


%%%%%%%%%%%%%%%% TRẢ LỜI NGẮN %%%%%%%%%%%%%%%%

% nguồn: Phan2_Lop_12_Dai_so_Chuong_1
\begin{bt}%[2D1C2-5]%[ma:4cd90faf]
	Đồ thị hàm số $y = x^3 - 2mx^2 + m^2x +  n$ có điểm cực tiểu là $I(1;3)$. Khi đó $m + n$ bằng
	\shortans{4} 
	 \loigiai{
		Ta có: $y = x^3 - 2mx^2 +  m^2x + n \Rightarrow y' = 3x^2 - 4mx + m^2 \Rightarrow y''=6x - 4m$\\
		Do $I(1;3)$ là điểm cực tiểu của đồ thị hàm số $y = x^3 - 2mx^2 +  m^2x + n$\\
		$\Rightarrow \heva{&I(1;2) \in y\\&y'(1) = 0\\&y''(1) > 0} \Leftrightarrow \heva{&1 - 2m + m^2 +  n = 3\\&m^2 - 4m + 3 = 0\\&6 - 4m > 0} \Leftrightarrow \heva{&1 -2m + m^2 + n = 3\\&\hoac{&m= 1\\&m = 3}\\&m < \dfrac{3}{2}} \Leftrightarrow \heva{&m = 1\\&n = 3}$\\
		Vậy $m + n  = 4$.
	} 
\end{bt}
